%! Unit Opener: Functions as Models — Unit 2, Lesson 2.0 — Slides (Beamer)
\documentclass[aspectratio=169, 11pt]{beamer}
\usepackage{atda-beamer}   % provides \royalheader{} and \sectionlabel[color]{}
\usepackage{pgfplots}      % atda-beamer does NOT load pgfplots
\pgfplotsset{compat=1.18}

\begin{document}

% TITLE SLIDE (CourseName is not defined in beamer, so the name is literal)
{
\setbeamercolor{background canvas}{bg=royal}
\begin{frame}[plain]
  \vspace{2.8cm}\hspace{0.55cm}%
  \begin{minipage}{0.9\textwidth}
    {\color{goldacc}\bfseries\small LESSON 2.0}\\[0.5cm]
    {\color{white}\bfseries\LARGE Unit Opener: Functions as Models}\\[1.2cm]
    {\color{mistmid}\small Algebra, Trigonometry, and Data Analysis~$\cdot$~Unit 2}
  \end{minipage}
\end{frame}
}

% ── HOOK ──────────────────────────────────────────────────────────────────────
\begin{frame}[t, plain]
  \royalheader{The Saturday Hike}
  \sectionlabel[goldacc]{hook}
  A friend hikes a mountain trail. Their watch records one thing --- elevation, every minute, for
  six hours --- and hands you a \textbf{picture}. No formula. No table.
  \begin{block}{Four things I want to know}
    How long did they hike? \quad How high did they get?\\
    When did they stop for lunch? \quad Did they come down faster than they went up?\\[2pt]
    \textbf{Can a picture answer all four?}
  \end{block}
\end{frame}

% ── THE GRAPH ─────────────────────────────────────────────────────────────────
\begin{frame}[t, plain]
  \royalheader{One Graph, Every Question}
  \sectionlabel{read it, do not solve it}
  \begin{center}
  \begin{tikzpicture}
  \begin{axis}[
      width=10.4cm, height=4.9cm,
      xlabel={$t$ (hours)}, ylabel={$E$ (feet)},
      xmin=0, xmax=6.4, ymin=0, ymax=2300,
      xtick={0,1,2,3,4,5,6}, ytick={0,800,1600,2000},
      grid=both, grid style={linegray!45}, axis lines=left,
      tick label style={font=\tiny}, label style={font=\tiny}]
    \addplot[royal, very thick, mark=*, mark size=1.4pt]
      coordinates {(0,800) (2,1600) (3,1600) (4,2000) (6,1000)};
  \end{axis}
  \end{tikzpicture}
  \end{center}
  \vspace{-0.2cm}
  Domain $0 \le t \le 6$ h \quad$\cdot$\quad starts at $800$ ft \quad$\cdot$\quad flat from
  $t=2$ to $3$ \quad$\cdot$\quad max $2000$ ft \textbf{at $t=4$}
\end{frame}

% ── THE RULE OF THE UNIT ──────────────────────────────────────────────────────
\begin{frame}[t, plain]
  \royalheader{The Rule of This Unit}
  \sectionlabel[goldacc]{why it matters}
  \begin{itemize}
    \setlength{\itemsep}{0.35cm}
    \item In this course a function's characteristics are read \textbf{from its graph}. That is
          the standard, not a shortcut.
    \item Every extreme gets reported two ways: \textbf{the value and where it happens}.
          ``$2000$ feet'' is half an answer.
    \item A graph may have \textbf{no zeros}. The hikers never reach sea level, so the curve
          never touches the $t$-axis.
    \item A \textbf{context sets the domain}. The model is in force for six hours and not one
          minute longer.
  \end{itemize}
\end{frame}

% ── THREE FAMILIES ────────────────────────────────────────────────────────────
\begin{frame}[t, plain]
  \royalheader{Three Families, Three Shapes}
  \sectionlabel{the parent functions}
  \begin{center}
  \begin{tabular}{@{}c@{\hspace{0.9cm}}c@{\hspace{0.9cm}}c@{}}
  \begin{tikzpicture}
  \begin{axis}[width=4.4cm, height=3.9cm, xmin=-3.2, xmax=3.2, ymin=-3.2, ymax=3.2,
      axis lines=middle, xtick=\empty, ytick=\empty]
    \addplot[royal, very thick, domain=-3:3, samples=2]{x};
  \end{axis}
  \end{tikzpicture}
  &
  \begin{tikzpicture}
  \begin{axis}[width=4.4cm, height=3.9cm, xmin=-3.2, xmax=3.2, ymin=-1.2, ymax=5.2,
      axis lines=middle, xtick=\empty, ytick=\empty]
    \addplot[royal, very thick, domain=-2.2:2.2, samples=60]{x^2};
  \end{axis}
  \end{tikzpicture}
  &
  \begin{tikzpicture}
  \begin{axis}[width=4.4cm, height=3.9cm, xmin=-3.2, xmax=3.2, ymin=-1.2, ymax=5.2,
      axis lines=middle, xtick=\empty, ytick=\empty]
    \addplot[royal, very thick, domain=-3:2.2, samples=60]{2^x};
  \end{axis}
  \end{tikzpicture}
  \\[-2pt]
  $f(x)=x$ & $g(x)=x^2$ & $h(x)=2^x$ \\
  linear & quadratic & exponential \\
  \end{tabular}
  \end{center}
\end{frame}

% ── HOW TO TELL THEM APART ────────────────────────────────────────────────────
\begin{frame}[t, plain]
  \royalheader{Adding or Multiplying?}
  \sectionlabel{telling the families apart}
  \renewcommand{\arraystretch}{1.4}
  \begin{tabularx}{\textwidth}{@{}p{3.0cm} p{4.6cm} X@{}}
    \textbf{Family} & \textbf{On a graph} & \textbf{In a table} \\
    \hline
    Linear & a straight line & adds the same amount each step \\
    Quadratic & one turning point, symmetric & the \emph{differences} change by the same amount \\
    Exponential & flattens toward a line it never touches & \textbf{multiplies} by the same
      factor each step \\
  \end{tabularx}
  \vspace{0.4cm}

  A line and an exponential both rise forever. \textbf{Adding versus multiplying is what
  separates them.}
\end{frame}

% ── ROADMAP ───────────────────────────────────────────────────────────────────
\begin{frame}[t, plain]
  \royalheader{The Unit 2 Roadmap}
  \sectionlabel{where we are going}
  \renewcommand{\arraystretch}{1.2}
  \small
  \begin{tabularx}{\textwidth}{@{}c p{2.9cm} X@{}}
    \textbf{Lesson} & \textbf{Standard} & \textbf{What you will be able to do} \\
    \hline
    2.1 & \texttt{AF.2a, e}      & Function notation; domain and range in context \\
    2.2 & \texttt{AF.1a--b}      & Parent functions and transformations \\
    2.3 & \texttt{AF.1d, f}      & Equation $\leftrightarrow$ graph, checked on the TI-84 \\
    2.4 & \texttt{AF.2b--d}      & Intercepts, zeros, extremes, intervals \\
    2.5 & \texttt{AF.2f--g}      & End behavior and asymptotes \\
    2.6 & \texttt{AF.2}          & Piecewise-defined functions \\
    2.7 & \texttt{AF.1c, e, g}, \texttt{AF.2h} & Choosing and comparing models \\
  \end{tabularx}
\end{frame}

% ── WHERE THIS LANGUAGE GETS SPENT ────────────────────────────────────────────
\begin{frame}[t, plain]
  \royalheader{Where This Language Gets Spent}
  \sectionlabel[goldacc]{the bridge}
  \begin{itemize}
    \setlength{\itemsep}{0.4cm}
    \item \textbf{Unit 1 is still underneath.} An excluded value becomes a \textbf{domain
          restriction} in 2.1 --- and in 2.5 it becomes a \textbf{vertical asymptote}.
    \item \textbf{Unit 3} takes the exponential family and asks how to undo it: logarithms, read
          with today's words.
    \item \textbf{Unit 6}'s normal curve and \textbf{Unit 8}'s sine wave are shapes you will read
          with these same five questions.
  \end{itemize}
\end{frame}

% ── ACTIVITY LAUNCH ───────────────────────────────────────────────────────────
\begin{frame}[t, plain]
  \royalheader{Group Activity: Reading the Day off a Graph}
  \sectionlabel{groups of three --- everyone starts at Tier R}
  Nobody hands you an equation today. You get pictures, and every answer is in the picture.
  \begin{block}{Your job}
    A rideshare fare, then a phone battery over ten hours: domain, range, every interval, and
    \textbf{both} absolute extremes --- value \emph{and} location. Then convert the clock times.
  \end{block}
  \vspace{0.2cm}
  Tier E: two students got it wrong. Find out how.
\end{frame}

% ── CLOSE ─────────────────────────────────────────────────────────────────────
\begin{frame}[t, plain]
  \royalheader{Exit Ticket}
  \sectionlabel[goldacc]{notes closed --- 5 minutes}
  A soccer ball is kicked from the ground; you get the graph of its height.
  \begin{enumerate}
    \setlength{\itemsep}{0.35cm}
    \item Domain and range, with units.
    \item The absolute maximum: value, location, and what it means.
    \item ``The domain of every function is all real numbers.'' Do you agree?
  \end{enumerate}
  \vspace{0.3cm}
  {\color{royal}\bfseries Next: Lesson 2.1 --- Function Notation, Domain, and Range in Context.}
\end{frame}

\end{document}
